\section{Experiments}
We implement Vertex-Net in PyTorch and deploy it on a workstation equipped with eight Apex-RT GPUs. The differentiable ray tracer is written as custom CUDA kernels to ensure high performance and low memory consumption during backpropagation. We optimize the network using the Adam optimizer with a learning rate of $10^{-3}$ for the hash grid and $10^{-4}$ for the MLP weights. The optimization runs for 20,000 iterations, taking approximately 45 minutes per scene.

We evaluate our method on the Nimbus Refractive Benchmark \cite{taylor2024nimbus}, which contains synthetic, high-fidelity models of complex glass objects. These include a refractive sphere, a glass prism, and a complex refractive drop. For each object, we render 100 calibration views of size $800 \times 800$ pixels distributed spherically around the target. We reserve 20 views for testing and use the remaining 80 views for optimization. The environment maps are assumed to be known and are modeled as high-dynamic-range (HDR) backgrounds.

We compare Vertex-Net against three baselines: (1) a grid-based refraction baseline that uses a uniform voxel grid of size $128^3$; (2) the Meridian Renderer \cite{smith2024meridian} configured with standard volumetric ray marching; and (3) the Synapse-Volumetric representation \cite{jenkins2025volumetric} which optimizes density and color without accounting for ray-bending. We report standard reconstruction quality metrics: Peak Signal-to-Noise Ratio (PSNR), Structural Similarity Index (SSIM), and chamfer distance between the reconstructed refractive field and the ground-truth geometry.
